\chapter{常微分方程解析解法}

\intro{
	任何一个能够作的比他的实际工作更为精巧的计算者，都决不会在数值计算上浪费时间。
	
	\rightline{——Carl Friedrich Gauss}
}

\section{一阶常微分方程的解析方法}

\subsection{分离变量法}

\begin{theorem}{分离变量法}
	若一阶微分方程可写为
	\[
	\frac{dy}{dt} = g(t) \, h(y)
	\]
	其中 $g(t)$ 仅为 $t$ 的函数，$h(y)$ 仅为 $y$ 的函数，且 $h(y) \neq 0$，则可通过分离变量：
	\[
	\frac{dy}{h(y)} = g(t) \, dt
	\]
	两边积分得隐式解：
	\[
	\int \frac{dy}{h(y)} = \int g(t) \, dt + C
	\]
\end{theorem}

\begin{example}{指数衰减方程}
	求解 $\displaystyle \frac{dy}{dt} = -a y, \; y(0) = y_0$。
\end{example}

\begin{solution}
	分离变量：$\displaystyle \frac{dy}{y} = -a \, dt$，积分得 $\ln |y| = -a t + C$，即 $y = C_1 e^{-a t}$。由初值条件得 $C_1 = y_0$，故
	\[
	y(t) = y_0 e^{-a t}
	\]
\end{solution}

\subsection{一阶线性微分方程与积分因子法}

一阶线性微分方程的标准形式为：
\begin{myformula}
	y' + p(t) y = q(t)
\end{myformula}

\begin{theorem}{积分因子法}
	一阶线性 ODE $y' + p(t) y = q(t)$ 的通解为：
	\[
	y(t) = e^{-\int p(t) dt} \left[ \int q(t) e^{\int p(t) dt} dt + C \right]
	\]
	其中 $\mu(t) = e^{\int p(t) dt}$ 称为\textbf{积分因子}。
\end{theorem}

\begin{proof}
	将方程两端乘以积分因子 $\mu(t) = e^{\int p(t) dt}$：
	\[
	e^{\int p(t) dt} y' + p(t) e^{\int p(t) dt} y = q(t) e^{\int p(t) dt}
	\]
	注意到左端为乘积导数：
	\[
	\frac{d}{dt}\left[ y \, e^{\int p(t) dt} \right] = q(t) e^{\int p(t) dt}
	\]
	积分即得结论。
\end{proof}

\section{二阶线性常微分方程}

\subsection{常系数齐次方程}

二阶常系数线性齐次微分方程的形式为：
\begin{myformula}
	a y'' + b y' + c y = 0, \quad a \neq 0
\end{myformula}

设解具有形式 $y = e^{\lambda t}$，代入得\textbf{特征方程}：
\begin{myformula}
	a \lambda^2 + b \lambda + c = 0
\end{myformula}

其判别式 $\Delta = b^2 - 4ac$。根据判别式的符号，通解有三种情形：

\begin{myTheorem}{二阶常系数齐次方程的通解结构}
	设特征根为 $\lambda_1, \lambda_2$，则通解为：
	\begin{enumerate}
		\item 若 $\Delta > 0$（两相异实根）：
		\[
		y(t) = C_1 e^{\lambda_1 t} + C_2 e^{\lambda_2 t}
		\]
		\item 若 $\Delta = 0$（重根 $\lambda = -b/(2a)$）：
		\[
		y(t) = (C_1 + C_2 t) e^{\lambda t}
		\]
		\item 若 $\Delta < 0$（共轭复根 $\lambda_{1,2} = \alpha \pm i\beta$）：
		\[
		y(t) = e^{\alpha t} (C_1 \cos \beta t + C_2 \sin \beta t)
		\]
	\end{enumerate}
\end{myTheorem}

\begin{proof}
	\textbf{重根情形}的推导：设 $\lambda$ 为二重根，则除 $y_1 = e^{\lambda t}$ 外，还需一个与 $y_1$ 线性无关的解。用常数变易法：令 $y_2(t) = v(t) e^{\lambda t}$，代入方程：
	\[
	a(v'' e^{\lambda t} + 2\lambda v' e^{\lambda t} + \lambda^2 v e^{\lambda t}) + b(v' e^{\lambda t} + \lambda v e^{\lambda t}) + c v e^{\lambda t} = 0
	\]
	整理并以 $a\lambda^2 + b\lambda + c = 0$ 及 $2a\lambda + b = 0$ 化简，得 $a v'' e^{\lambda t} = 0$，故 $v'' = 0$，$v(t) = C_1 + C_2 t$。取 $v(t) = t$，得第二个线性无关解 $y_2 = t e^{\lambda t}$。
\end{proof}

\subsection{常系数非齐次方程}

非齐次方程的形式为：
\begin{myformula}
	a y'' + b y' + c y = f(t)
\end{myformula}

\begin{theorem}{非齐次方程的通解结构}
	设 $y_h(t)$ 为对应齐次方程的通解，$y_p(t)$ 为非齐次方程的一个特解，则非齐次方程的通解为：
	\[
	y(t) = y_h(t) + y_p(t)
	\]
\end{theorem}

\textbf{求特解的常用方法：}

\medskip
\textbf{1. 待定系数法}——适用于 $f(t)$ 为多项式、指数函数、正弦/余弦函数或其乘积的情形。

\begin{example}{指数型驱动力}
	求解 $y'' + 3y' + 2y = e^{-t}$。
\end{example}

\begin{solution}
	特征方程 $\lambda^2 + 3\lambda + 2 = 0$，根为 $\lambda_1 = -1, \lambda_2 = -2$。齐次通解 $y_h = C_1 e^{-t} + C_2 e^{-2t}$。
	
	由于 $-1$ 为单重特征根，设特解 $y_p = A t e^{-t}$。求导后代回方程得 $A = 1$。故通解：
	\[
	y(t) = C_1 e^{-t} + C_2 e^{-2t} + t e^{-t}
	\]
\end{solution}

\medskip
\textbf{2. 拉普拉斯变换法}——适用于具有初始条件的初值问题。

\medskip
\textbf{3. 常数变易法}——适用于一般的 $f(t)$：

\begin{theorem}{常数变易公式}
	设齐次方程的基本解组为 $y_1(t), y_2(t)$，其 Wronsky 行列式为 $W(t)$，则特解为：
	\[
	y_p(t) = -y_1(t) \int \frac{y_2(t) f(t)}{a W(t)} \, dt + y_2(t) \int \frac{y_1(t) f(t)}{a W(t)} \, dt
	\]
\end{theorem}

\section{一阶线性常微分方程组}

一阶线性方程组的标准形式：
\begin{myformula}
	\mathbf{y}' = A(t) \, \mathbf{y} + \mathbf{g}(t)
\end{myformula}
其中 $\mathbf{y} \in \mathbb{R}^n$，$A(t) \in \mathbb{R}^{n \times n}$。

\subsection{常系数齐次方程组}

若 $A$ 为常矩阵，齐次系统 $\mathbf{y}' = A\mathbf{y}$ 的通解为：
\begin{myformula}
	\mathbf{y}(t) = e^{At} \, \mathbf{y}_0
\end{myformula}

其中矩阵指数定义为：
\begin{myformula}
	e^{At} = \sum_{k=0}^{\infty} \frac{(At)^k}{k!}
\end{myformula}

\begin{theorem}{矩阵指数的计算方法（对角化）}
	若 $A$ 可对角化，即存在可逆矩阵 $P$ 使 $A = P \Lambda P^{-1}$，其中 $\Lambda = \operatorname{diag}(\lambda_1, \ldots, \lambda_n)$，则：
	\[
	e^{At} = P \, \operatorname{diag}(e^{\lambda_1 t}, \ldots, e^{\lambda_n t}) \, P^{-1}
	\]
\end{theorem}

\begin{proof}
	由定义：
	\[
	e^{At} = \sum_{k=0}^{\infty} \frac{(P\Lambda P^{-1})^k t^k}{k!} = P \left( \sum_{k=0}^{\infty} \frac{\Lambda^k t^k}{k!} \right) P^{-1} = P e^{\Lambda t} P^{-1}
	\]
	而 $\Lambda^k = \operatorname{diag}(\lambda_1^k, \ldots, \lambda_n^k)$，故 $e^{\Lambda t} = \operatorname{diag}(e^{\lambda_1 t}, \ldots, e^{\lambda_n t})$。
\end{proof}

\section{动力系统定性分析初步}

\subsection{平衡点与稳定性}

\begin{mydefinition}{平衡点（不动点）}
	对于自治系统 $\dot{\mathbf{y}} = \mathbf{f}(\mathbf{y})$，若 $\mathbf{f}(\mathbf{y}^*) = 0$，则称 $\mathbf{y}^*$ 为系统的平衡点。
\end{mydefinition}

\begin{mydefinition}{Lyapunov 稳定性}
	平衡点 $\mathbf{y}^*$ 称为
	\begin{enumerate}
		\item \textbf{稳定的}：$\forall \varepsilon > 0, \exists \delta > 0$，使 $\|\mathbf{y}(0) - \mathbf{y}^*\| < \delta$ 时，$\|\mathbf{y}(t) - \mathbf{y}^*\| < \varepsilon$ 对所有 $t \ge 0$ 成立；
		\item \textbf{渐近稳定的}：稳定且 $\lim_{t\to\infty} \|\mathbf{y}(t) - \mathbf{y}^*\| = 0$；
		\item \textbf{不稳定的}：不满足稳定性条件。
	\end{enumerate}
\end{mydefinition}

\subsection{二维线性系统的相图分类}

对于系统 $\dot{\mathbf{y}} = A\mathbf{y}$（$A \in \mathbb{R}^{2\times2}$），根据 $A$ 的特征值可将平衡点分类：

\begin{table}[H]
	\centering
	\caption{二维线性系统平衡点分类}
	\begin{tabular}{ccc}
		\toprule
		特征值类型 & 平衡点类型 & 稳定性 \\
		\midrule
		$\lambda_1 < \lambda_2 < 0$ & 稳定结点 & 渐近稳定 \\
		$\lambda_1 > \lambda_2 > 0$ & 不稳定结点 & 不稳定 \\
		$\lambda_1 < 0 < \lambda_2$ & 鞍点 & 不稳定 \\
		$\lambda_{1,2} = \alpha \pm i\beta,\; \alpha < 0$ & 稳定焦点 & 渐近稳定 \\
		$\lambda_{1,2} = \alpha \pm i\beta,\; \alpha > 0$ & 不稳定焦点 & 不稳定 \\
		$\lambda_{1,2} = \pm i\beta$ & 中心 & 稳定 \\
		\bottomrule
	\end{tabular}
\end{table}

\begin{remark}
	对于非线性系统 $\dot{\mathbf{y}} = \mathbf{f}(\mathbf{y})$，可在平衡点附近进行线性化：令 $\mathbf{f}(\mathbf{y}) \approx J(\mathbf{y}^*)(\mathbf{y} - \mathbf{y}^*)$，其中 $J$ 为 Jacobi 矩阵。若线性化系统的平衡点为渐近稳定的（所有特征值的实部均为负），则原非线性系统的平衡点也是渐近稳定的。
\end{remark}
